% ---------------------------------------------------------------
% Abkürzungs- und Symbolverzeichnis (Vorspann)
% ---------------------------------------------------------------
\addchap{Abkürzungsverzeichnis}
\begingroup
\renewcommand{\theHtable}{abkverz.\arabic{table}}
\small
\renewcommand{\arraystretch}{1.12}
\begin{longtable}{@{}p{2.3cm}p{12.4cm}@{}}
A3C & Asynchronous Advantage Actor-Critic \\
AC & Actor-Critic \\
ACK & Acknowledgement (TCP-Flag) \\
AE-RL & Adversarial Environment Reinforcement Learning \\
AWGN & Additive White Gaussian Noise, additives weißes gaußsches Rauschen \\
BIM & Basic Iterative Method \\
C-LSTM & kombinierte CNN-LSTM-Architektur \\
CNN & Convolutional Neural Network \\
CSV & Comma-Separated Values \\
CUDA & Compute Unified Device Architecture \\
DDoS & Distributed Denial of Service \\
DDPG & Deep Deterministic Policy Gradient \\
DDQN & Double Deep Q-Network \\
DQN & Deep Q-Network \\
DRL & Deep Reinforcement Learning \\
ELA & Equalized Loss of Accuracy \\
FF & Forschungsfrage \\
FGSM & Fast Gradient Sign Method \\
FIN & Finish (TCP-Flag) \\
FN & False Negative, verpasster Angriff \\
FP & False Positive, Fehlalarm \\
GPU & Graphics Processing Unit \\
HTTP & Hypertext Transfer Protocol \\
IDS & Intrusion Detection System \\
IoT & Internet of Things \\
IP & Internet Protocol \\
KI & Künstliche Intelligenz \\
$k$-NN & $k$-Nearest-Neighbors \\
LSTM & Long Short-Term Memory \\
MDP & Markov-Entscheidungsprozess (Markov Decision Process) \\
ML & maschinelles Lernen (Machine Learning) \\
MLP & Multilayer Perceptron \\
MSE & mittlerer quadratischer Fehler (Mean Squared Error) \\
NIDS & Network Intrusion Detection System \\
NN & neuronales Netz \\
PCAP & Packet Capture (Mitschnittformat) \\
PG & Policy Gradient \\
PPO & Proximal Policy Optimization \\
RBF & radiale Basisfunktion \\
ReLU & Rectified Linear Unit \\
RF & Random Forest \\
RNG & Zufallszahlengenerator (Random Number Generator) \\
RL & Reinforcement Learning \\
RLA & Relative Loss of Accuracy \\
RNN & rekurrentes neuronales Netz \\
RST & Reset (TCP-Flag) \\
SA-MDP & State-Adversarial Markov Decision Process \\
SDN & Software-Defined Networking \\
SMOTE & Synthetic Minority Over-sampling Technique \\
SNR & Signal-Rausch-Verhältnis (Signal-to-Noise Ratio) \\
SVM & Support Vector Machine \\
SYN & Synchronize (TCP-Flag) \\
TCP & Transmission Control Protocol \\
TN & True Negative \\
TP & True Positive \\
UDP & User Datagram Protocol \\
XAI & Explainable AI \\
\end{longtable}
\endgroup

\addchap{Symbolverzeichnis}
\begingroup
\renewcommand{\theHtable}{symverz.\arabic{table}}
\small
\renewcommand{\arraystretch}{1.18}
\renewcommand{\arraystretch}{1.06}
\begin{longtable}{@{}p{3.0cm}p{11.7cm}@{}}
\multicolumn{2}{@{}l}{\textbf{Generatormodell und Kanal}}\\[2pt]
$\mathbf{x}\in\mathbb{R}^{8}$ & unverrauschter Merkmalsvektor eines Flusses \\
$\mathbf{y}\in\mathbb{R}^{8}$ & beobachteter, verrauschter Merkmalsvektor \\
$\mathbf{n}$ & Rauschvektor des AWGN-Kanals \\
$c$ & Klassenlabel, $0$ Normal, $1$ DDoS \\
$\pi$ & Angriffsanteil (Prior der Klasse DDoS), $\pi=0{,}3$ \\
$\boldsymbol{\mu}_c$, $\Delta\mu_j$ & Klassenmittelwerte, Mittelwertabstand im Merkmal $j$ \\
$\sigma$ & intrinsische Klassenstreuung je Merkmal, $\sigma=0{,}20$ \\
$P_j$ & Signalleistung des Merkmals $j$ \\
$\sigma_{n,j}^2$, $\boldsymbol{\Sigma}_n$ & Rauschvarianz des Merkmals $j$, Kovarianzmatrix des Rauschens \\
$h$ & Kanalfaktor, hier $h=1$ \\[4pt]
\multicolumn{2}{@{}l}{\textbf{Entscheidungstheorie}}\\[2pt]
$p$ & Angriffs-Posterior $P(\mathrm{DDoS}\mid\mathbf{y})$ \\
$t$ & Posterior-Schwelle eines Modells, implizite Schwelle \\
$D$, $D^2$ & Mahalanobis-Abstand der Klassenmittel \\
$\Lambda(\mathbf{y})$ & Log-Likelihood-Verhältnis (Prüfgröße) \\
$\mathbf{w}$ & Normalenvektor der Bayes-Grenze, $w_j=\Delta\mu_j/(\sigma^2+\sigma_{n,j}^2)$ \\
$\lambda$ & Schwellenverschiebung durch den Prior, $\lambda=\ln\bigl((1-\pi)/\pi\bigr)$ \\
$\tau(t)$ & Schwelle auf $\Lambda$, die der Posterior-Schwelle $t$ entspricht \\
$P_e$, $P_e^{\pi}$ & Bayes-Fehlerwahrscheinlichkeit bei gleichen bzw. tatsächlichen Priors \\
$Q_{\mathrm{G}}(z)$ & Gaußsche Q-Funktion \\[4pt]
\multicolumn{2}{@{}l}{\textbf{Reinforcement Learning und Training}}\\[2pt]
$s$, $s'$ & Zustand und Folgezustand im allgemeinen MDP \\
$a$ & Aktion, hier die Entscheidung Normal ($0$) oder DDoS ($1$) \\
$r$, $G_k$ & Reward, diskontierter Return ab Schritt $k$ \\
$R_{\mathrm{TP}},\dots,R_{\mathrm{FN}}$ & Einträge der Reward-Matrix \\
$\gamma$ & Diskontierungsfaktor \\
$d$ & Terminal-Flag einer Transition \\
$Q(\mathbf{y},a)$ & Aktionswertfunktion (Q-Wert) \\
$Q_{\mathrm{H}}$ & Q-Wert bei Regression mit Huber-Verlust \\
$\theta$, $\bar{\theta}$ & Parameter von Online- und Zielnetz \\
$Y$ & Regressionsziel des Q-Netzes \\
$\ell_{\mathrm{H}}$ & Huber-Verlust ($\beta=1$) \\
$\epsilon$, $\alpha$ & Explorationsrate, Lernrate \\
$\eta^2$, $\Omega$ & Varianz des Eingangsrauschens und Tikhonov-Strafterm nach Bishop \\[4pt]
\multicolumn{2}{@{}l}{\textbf{Bewertung}}\\[2pt]
$\mathrm{F1}_s$ & F1-Score bei der SNR-Stufe $s$ \\
$\mathrm{RLA}_s$, $\mathrm{ELA}_s$ & Relative und Equalized Loss of Accuracy nach Sáez et al. \\
\end{longtable}
\endgroup
